% Options for packages loaded elsewhere
\PassOptionsToPackage{unicode}{hyperref}
\PassOptionsToPackage{hyphens}{url}
\documentclass[
]{article}
\usepackage{xcolor}
\usepackage{amsmath,amssymb}
\setcounter{secnumdepth}{-\maxdimen} % remove section numbering
\usepackage{iftex}
\ifPDFTeX
  \usepackage[T1]{fontenc}
  \usepackage[utf8]{inputenc}
  \usepackage{textcomp} % provide euro and other symbols
\else % if luatex or xetex
  \usepackage{unicode-math} % this also loads fontspec
  \defaultfontfeatures{Scale=MatchLowercase}
  \defaultfontfeatures[\rmfamily]{Ligatures=TeX,Scale=1}
\fi
\usepackage{lmodern}
\ifPDFTeX\else
  % xetex/luatex font selection
\fi
% Use upquote if available, for straight quotes in verbatim environments
\IfFileExists{upquote.sty}{\usepackage{upquote}}{}
\IfFileExists{microtype.sty}{% use microtype if available
  \usepackage[]{microtype}
  \UseMicrotypeSet[protrusion]{basicmath} % disable protrusion for tt fonts
}{}
\makeatletter
\@ifundefined{KOMAClassName}{% if non-KOMA class
  \IfFileExists{parskip.sty}{%
    \usepackage{parskip}
  }{% else
    \setlength{\parindent}{0pt}
    \setlength{\parskip}{6pt plus 2pt minus 1pt}}
}{% if KOMA class
  \KOMAoptions{parskip=half}}
\makeatother
\setlength{\emergencystretch}{3em} % prevent overfull lines
\providecommand{\tightlist}{%
  \setlength{\itemsep}{0pt}\setlength{\parskip}{0pt}}
\usepackage{bookmark}
\IfFileExists{xurl.sty}{\usepackage{xurl}}{} % add URL line breaks if available
\urlstyle{same}
\hypersetup{
  pdftitle={20260223 Anotações sobre The Eye of the Master},
  hidelinks,
  pdfcreator={LaTeX via pandoc}}

\title{20260223 Anotações sobre The Eye of the Master}
\author{}
\date{2026-02-23}

\begin{document}
\maketitle

up:: 099 MOC Anotações

7

\section{A invenção de máquinas depende do conhecimento da
função/trabalho que ela
substitui}\label{a-invenuxe7uxe3o-de-muxe1quinas-depende-do-conhecimento-da-funuxe7uxe3otrabalho-que-ela-substitui}

\begin{quote}
"When industrial machines such as looms and lathes were invented, in
fact, it was not thanks to the solitary genius of an engineer but
through the imitation of the collective diagram of labour: \emph{by
capturing the patterns of hand movements and tools, the subdued
creativity of workers' know-how, and turning them into mechanical
artefacts}. Following this theory of invention, which was already shared
by Smith, Babbage, and Marx in the nineteenth century, this book {[}The
Eye of the Master{]} argues that the most sophisticated `intelligent'
machines have also emerged by imitating the outline of the collective
division of labour." (Pasquinelli, 2023, p. 6; grifo meu)
\end{quote}

Alan Turing, falando sobre a \emph{Automatic Computing Engine} (ACE),
divide seus operadores em "mestres" e "servos":

\begin{quote}
"Its masters will plan out instruction tables for it, thinking up deeper
and deeper ways of using it. Its servants will feed it with cards as it
calls for them. They will put right any parts that go wrong. They will
assemble data that it requires. In fact the servants will take the place
of limbs." (\emph{apud} Pasquinelli, p. 7)
\end{quote}

Contudo, enquanto o trabalho dos servos é passível de substituição,
também o trabalho \emph{dos mestres} pode sê-lo:

\begin{quote}
"The masters are liable to get replaced because{[},{]} as soon as any
technique becomes at all stereotyped{[},{]} it becomes possible to
devise a system of instruction tables which will enable the electronic
computer to do it for itself. It may happen however that the masters
will refuse to do this. They may be unwilling to let their jobs be
stolen from them in this way. In that case{[},{]} \emph{they would
surround the whole of their work with mystery and make excuses, couched
in well chosen gibberish}, whenever any dangerous suggestions were
made." (Ibid.)
\end{quote}

É em sentido análogo que Marx fala que A maquinaria tem caráter
revolucionário, pois

\begin{quote}
"A grande indústria rasgou o véu que ocultava aos homens seu próprio
processo social de produção e que convertia os diversos ramos da
produção, que se haviam particularizado de modo natural-espontâneo, em
enigmas uns em relação aos outros, e inclusive para o iniciado em cada
um desses ramos. O princípio da grande indústria, a saber, o de
dissolver cada processo de produção propriamente dito em seus elementos
constitutivos, e, antes de tudo, fazê-lo sem nenhuma consideração para
com a mão humana, criou a mais moderna ciência da tecnologia." (Marx,
2017, p. 556)"
\end{quote}

Ao que conclui que "A indústria moderna jamais considera nem trata como
definitiva a forma existente de um processo de produção. Sua base
técnica é, por isso, revolucionária, ao passo que a de todos os modos de
produção anteriores era essencialmente conservadora." {[}@Marx2017 p.
557{]}

\section{Algoritmos como processos de
trabalho}\label{algoritmos-como-processos-de-trabalho}

\begin{quote}
"{[}Donald{]} Knuth observed {[}in his article \emph{Ancient Babylonian
Algorithms}, 1972{]} that mathematical formulas that today would be
defined as algebraic or analytical were already described by the
Babylonians through step-by-step procedures, namely algorithms. These
procedures were, of course, formulated in the words of the common
language and not yet in the symbolic metalanguage of mathematics.
Knuth's research confirms the hypothesis that \emph{procedure-based
methods} (what he called a `machine language') \emph{predated the
consolidation of mathematics as a metalanguage of symbolic
representations} {[}...{]}" {[}@Pasquinelli2022 p. 29; grifo meu{]}
\end{quote}

\begin{center}\rule{0.5\linewidth}{0.5pt}\end{center}

\subsubsection{Referências}\label{referuxeancias}

\end{document}
